% BHU PYQ -- Transcribed & Corrected
% Paper: PHYSE-11 / PHYSE11: Mathematical Problem Solving Skills
% Exam: B.Sc. (Hons.) Semester I Examination, 2024-25 (NEP)
% Category: Skill Enhancement Course (SEC)
% Department: Physics

\documentclass[11pt,a4paper]{article}
\usepackage[a4paper,top=2.5cm,bottom=2.5cm,left=2.8cm,right=2.8cm,headheight=14pt]{geometry}
\usepackage{amsmath,amssymb}
\usepackage[T1]{fontenc}
\usepackage[utf8]{inputenc}
\usepackage{lmodern,microtype,fancyhdr,enumitem,hyperref,lastpage,parskip}

\hypersetup{
    colorlinks=true,
    urlcolor=black,
    linkcolor=black,
    pdftitle={PHYSE11: Mathematical Problem Solving Skills -- BHU Sem I 2024-25 (NEP SEC)}
}

\pagestyle{fancy}
\fancyhf{}
\renewcommand{\headrulewidth}{0.4pt}
\renewcommand{\footrulewidth}{0.4pt}
\fancyhead[L]{\small   \textit{Banaras Hindu University}}
\fancyhead[R]{\small   \textit{PHYSE-11: Math Problem Solving Skills (SEC)}}
\fancyfoot[C]{\small Page  \thepage\ of \pageref{LastPage}}


\newlist{parts}{enumerate}{2}
\setlist[parts,1]{label=(\alph*),leftmargin=*,itemsep=4pt,topsep=2pt}
\setlist[parts,2]{label=(\roman*),leftmargin=*,itemsep=2pt,topsep=2pt}

\newcommand{\pts}[1]{\hfill   \textit{[#1]}}


\begin{document}


\begin{center}
  {\large\bfseries BANARAS HINDU UNIVERSITY}\\[2pt]
  {
\normalsize B.Sc. (Hons.) --- Semester I Examination, 2024-25 (NEP)}\\[6pt]
  \rule{0.8\linewidth}{0.5pt}\\[6pt]
  {\large\bfseries PHYSICS}\\[4pt]
  {\large\bfseries Paper Code: PHYSE-11 : Mathematical Problem Solving Skills}\\[4pt]
  {
\normalsize Time: 2:30 Hours\hfill Maximum Marks: 70}\\[2pt]
  {\small REF NO. CECONF/N/2024-25/1}
\end{center}

\rule{\linewidth}{0.4pt}

\smallskip



\noindent   \textit{Instructions: Answer   \textbf{five} questions in all including Question No.~1 which is compulsory. All questions carry equal marks (14 marks each).

\bigskip




\noindent   \textbf{Question 1.} Answer any four of the following:\pts{$3.5  \times 4 = 14$}

\begin{parts}
  \item Explain the term   \textit{integrating factor} used for solving a first-order differential equation.\pts{3.5}
  \item Find the roots of $\sin (2x) = 0$ in the interval $(-\pi, \pi)$.\pts{3.5}
  \item Prove that the vector field $\vec{F}(x, y, z) = y^2 \hat{i} + x^3 y \hat{j} - z \hat{k}$ is a rotational field.\pts{3.5}
  \item Find the divergence of the vector $\vec{A} = 4 \hat{i} - 9 \hat{j} - 9 \hat{k}$.\pts{3.5}
  \item What are dependent and independent variables? Explain with suitable physical examples.\pts{3.5}
  \item Write the Gaussian probability distribution formula and mention its importance in physics.\pts{3.5}
\end{parts}

\medskip\rule{\linewidth}{0.2pt}\medskip




\noindent   \textbf{Question 2.}

\begin{parts}
  \item What is an ordinary differential equation? Discuss the different analytical methods used for solving first-order differential equations.\pts{6}
  \item Solve the following first-order differential equation:\pts{8}
  \[
    x \frac{dy}{dx} + 3y = 4x^2 - 3x
  \]
\end{parts}

\medskip\rule{\linewidth}{0.2pt}\medskip




\noindent   \textbf{Question 3.}

\begin{parts}
  \item Explain the terms   \textit{divergence} and   \textit{curl} of vector fields. Discuss their physical significance in electromagnetism and fluid dynamics.\pts{7}
  \item Find the vector ($\vec{A}  \times \vec{B}$) and scalar ($\vec{A} \cdot \vec{B}$) products of the vectors:\pts{7}
  \[
    \vec{A} = 3\hat{i} + 7\hat{j} - 8\hat{k} \quad  \text{and} \quad \vec{B} = -2\hat{i} + 5\hat{j} + 3\hat{k}
  \]
\end{parts}

\medskip\rule{\linewidth}{0.2pt}\medskip




\noindent   \textbf{Question 4.} Find the degree 4 Taylor polynomial $P_4(x)$ for the function $f(x) = \sqrt{x}$ expanded about $x = 4$.\pts{14}

\medskip\rule{\linewidth}{0.2pt}\medskip




\noindent   \textbf{Question 5.} Plot and analyze the quadratic function $f(x) = x^2 - ax + bx + ab$ for the following cases:

\begin{parts}
  \item $a > 0$, $a = 0$, and $a < 0$ (with $b > 0$).\pts{7}
  \item Calculate the roots for $a > 0$ and $b < 0$.\pts{3}
  \item Using the second derivative test, find the local minima for $a > 0$ and $b < 0$.\pts{4}
\end{parts}

\medskip\rule{\linewidth}{0.2pt}\medskip




\noindent   \textbf{Question 6.}

\begin{parts}
  \item Consider the experiment of flipping a fair coin four times. Calculate the probability that exactly one head is obtained.\pts{3}
  \item Calculate the probability that the first three flips are tails.\pts{5}
  \item Show that if we throw three fair dice, the expected number of 1's that will appear is $0.5$.\pts{6}
\end{parts}

\medskip\rule{\linewidth}{0.2pt}\medskip




\noindent   \textbf{Question 7.} Consider the following marks obtained by 15 students in a class:
\[
  \{7, 7, 4, 3, 4, 4, 5, 4, 6, 4, 4, 7, 3, 3, 5\}
\]

\begin{parts}
  \item Convert the raw data into a discrete probability distribution.\pts{7}
  \item Calculate the mean ($\mu$) using the probability distribution.\pts{3}
  \item Calculate the standard deviation ($\sigma$) using the probability distribution.\pts{4}
\end{parts}

\vfill
\rule{\linewidth}{0.4pt}

\begin{center}\small
\href{https://bhu-exam.vercel.app/}{Click here to visit our website}\\[4pt]
\href{https://me-aryan.vercel.app/}{Click here to visit our contributor}
\end{center}

\end{document}
